\clearpage
\thispagestyle{plain}
\setcounter{page}{2}
\singlespacing

{\centering\bfseries РЕФЕРАТ\par}\vspace{0.6ex}

\noindent
Объём работы --- 91 с., 17 табл., 24 рис., 15 источников.

\textbf{Ключевые слова:} акселерометр, генератор случайных чисел, дробовой
шум, лавинный пробой, операционный усилитель, постобработка, тепловой шум,
энтропия, NIST STS, TRNG.

\vspace{0.4ex}

\textbf{Постановка задачи.} В работе ставится задача сравнительного анализа
физических источников энтропии и разработки на основе наиболее подходящего из
них концепции компактного физического генератора случайных чисел (ГИСП),
пригодного для практической реализации на доступной элементной базе. Полученная
система должна удовлетворять формальным критериям случайности по методике
пакета NIST STS, обеспечивать достаточную скорость генерации и предоставлять
выходной поток, пригодный для применения в криптографических задачах
(формирование начального состояния ГПСЧ, генерация ключей и т.\,п.).

\textbf{Методы решения.} Теоретический анализ восьми классов физических
источников энтропии (тепловой, дробовой, лавинный, квантовый, акустический,
электромагнитный и др.) с привязкой каждого к разделам математического
описания, экспериментальная регистрация шести каналов на стенде Arduino Nano
с двухкаскадным усилительным трактом на операционных усилителях LM358 / TL072,
постобработка по схеме фон Неймана и каскадных XOR-децимирующих алгоритмов,
формальная проверка случайности пакетом NIST Statistical Test Suite 2.1.2,
сравнительная оценка по совокупности критериев (PASS/FAIL NIST, минимальная
энтропия, скорость, сложность аппаратной реализации).

\textbf{Основные результаты.} Сконструирован и натурно проверен стенд,
реализующий шесть независимых источников энтропии. Установлено, что наилучший
профиль NIST STS обеспечивает канал лавинного шума обратносмещённого BE-перехода
биполярного транзистора BC337 в финальной конфигурации с двухступенчатой
постобработкой (фон Нейман $+$ XOR-децимация с окном 8): получено
\textbf{15 PASS / 0 FAIL / 0 N/A} на всех пятнадцати группах тестов NIST STS;
минимальная энтропия выходного потока $0{,}9999$ бит на бит, скорость генерации
после каскадной постобработки $\sim 1{,}6$~кбит/с при стоимости элементной базы
менее 50 рублей. Предложены направления дальнейшего развития: интеграция
второго аналогичного канала с XOR-смешением, переход на печатную плату,
упаковка в законченный модуль с собственным микроконтроллером.
